\section{Classical vs Quantum Communications}
Quando parliamo di trasferimento di informazione pensiamo a quelli che oggi sono gli strumenti utilizzati dal mondo classico (un esempio classico è internet, le comunicazioni seriali ecc.), ma quando parliamo del mondo quantistico la situazione è un po' diversa.

Nel mondo classico l'informazione viene codificata all'interno dei \textbf{bit} che poi vengono trasmessi. Mentre nel mondo quantum, l'informazione viene codificata all'interno dei \textbf{Qubit}. Di conseguenza si potrebbe pensare che la trasmissione di informazione quantistica sia solo legata alla trasmissione dei Qubit, ma questo è errato.

Di base uno strumento di comunicazione classico trasmette dei bit, mentre uno strumento di comunicazione quantistico sfrutta l'entanglement e la condivisione di bit classici. Quindi la comunicazione quantistica dipende ancora da quella classica. Questo crea quel collo di bottiglia per la velocità di trasmissione di informazioni. Ciò è imprescindibile, quindi non si potrà fare a meno di comunicazioni classiche per le trasmissioni quantistiche. 

Uno schema che può ben rappresentare come sono classificate le possibili reti è il seguente:

\begin{table}[h]
\centering
\begin{tabular}{|>{\centering\arraybackslash}p{3cm}|>{\centering\arraybackslash}p{3.5cm}|>{\centering\arraybackslash}p{3.5cm}|}
\hline
 & \textbf{Classical Communications} & \textbf{Quantum Communications} \\
\hline
\textbf{Classical Message (bits)} & \textbf{Internet} & \textbf{QKD \newline SuperDense Coding} \\
\hline
\textbf{Quantum Message (qubits)} & & \textbf{Quantum Internet} \\
\hline
\end{tabular}
\caption{Classical vs Quantum Communications}
\end{table}

Da tale tabella, quello che si comprende, è come vengano sfruttate le comunicazioni quantistiche. Il \textbf{QKD} ed il \textbf{Super Dense Coding} sono tecniche basate su entanglement ma per comunicazioni classiche (quindi condivisione di chiavi ecc.). Mentre nel caso della \textbf{Quantum Internet}, quello che si cerca di fare è il vero e proprio "trasferimento" di un qubit (non un clone, attenzione, ricordiamo il teorema del no-cloning).

Quindi alla fine le reti possono essere classificate in due principali aree:
\begin{itemize}
    \item \textbf{Unentangled Networks}: Utili solamente per la QKD (Quantum Key Distribution), che sfrutta l'entanglement solo per la grande segretezza della sua trasmissione (lo si vedrà più in avanti). L'applicazione di tali tecniche è più facile rispetto ad altre tecniche (anche se questo non voglia dire che siano facili da implementare)
    
    \item \textbf{Entangled Networks}: Utili per il trasferimento di qubit interi (non di copia, ma di "spostamento"). Rispetto alle reti Unentangled, le reti entangled sono molto più difficili da utilizzare (o costruire).
    
\end{itemize}